% Source: https://artofproblemsolving.com/community/c3h1923090_how_to_find_sum_of_all_the_terms_in_an_arithmetic_sequence_formulaexplanation
% Stable identity: aops-topic-1923090-q01
\textbf{(a)} Let the terms be $a,a+d,\dots,a+(n-1)d=\ell$, with sum $S$. Write the sum forwards and backwards:
\[
\begin{aligned}
S&=a+(a+d)+\cdots+(\ell-d)+\ell,\\
S&=\ell+(\ell-d)+\cdots+(a+d)+a .
\end{aligned}
\]
In each of the $n$ columns the two terms add to $a+\ell$ (going right, the top term increases by $d$ and the bottom one decreases by $d$). So $2S=n(a+\ell)$ and $S=\dfrac{a+\ell}{2}\cdot n$.

\textbf{(b)} The terms are $3+4k$ for $k=0,1,\dots,24$, so there are $\dfrac{99-3}{4}+1=25$ terms and the sum is $\dfrac{3+99}{2}\cdot25=\boxed{1275}$.

\textbf{(c)} For $1,2,4$ (not arithmetic), the formula gives $\dfrac{1+4}{2}\cdot3=7.5$, but the sum is $7$. The pairing in (a) needs the opposite terms to have equal sums, which holds for arithmetic sequences.
